% IEEE Transactions on Computer-Aided Design Template
% GLACIER-MAESTRO Paper

\documentclass[journal]{IEEEtran}

% Packages
\usepackage{amsmath,amssymb,amsfonts}
\usepackage{algorithmic}
\usepackage{algorithm}
\usepackage{array}
\usepackage{graphicx}
\usepackage{textcomp}
\usepackage{url}
\usepackage{cite}
\usepackage{booktabs}
\usepackage{multirow}
\usepackage{color}
\usepackage{balance}

% Correct bad hyphenation here
\hyphenation{op-tical net-works semi-conduc-tor}

\begin{document}

\title{GLACIER-MAESTRO: A Comprehensive Framework for Robust Nonlinear Circuit Simulation Combining Logarithmic Transformation with Topology-Aware Strategy Orchestration}

\author{[Author Names]%
\thanks{Manuscript received [date]; revised [date]. This work was supported in part by [funding].}%
\thanks{[Authors] are with [Institution] (e-mail: [email]).}%
\thanks{Digital Object Identifier [DOI]}}

% The paper headers
\markboth{IEEE Transactions on Computer-Aided Design of Integrated Circuits and Systems, Vol.~XX, No.~XX, Month~2025}%
{Author \MakeLowercase{\textit{et al.}}: GLACIER-MAESTRO: Robust Nonlinear Circuit Simulation}

% Make the title area
\maketitle

\begin{abstract}
This paper presents GLACIER-MAESTRO, a revolutionary circuit simulation framework that achieves 100\% convergence on previously unsolvable nonlinear circuits. The framework combines two synergistic innovations: GLACIER (Gradient Logarithmic Adaptive Circuit Intelligent Exploration Resolver), which introduces logarithmic transformation directly into the Newton-Raphson loop with gradient-aware region identification, and MAESTRO (Multi-strategy Adaptive Engine for Smart Topology-driven Resolution and Orchestration), which exploits circuit topology to intelligently select and apply solving strategies. 

We demonstrate that GLACIER alone achieves convergence on circuits with LED saturation currents as low as $10^{-38}$ A, while MAESTRO's progressive activation and topology-aware strategies achieve 92.3\% success rate on a comprehensive benchmark of 52 challenging circuits. When combined, the framework achieves perfect 100\% convergence across all test cases, compared to only 36.5\% for traditional Newton-Raphson methods.

The key contributions include: (1) Phase 0 gradient-aware region identification that detects sharp nonlinearities before solving begins, (2) Full logarithmic transformation integration with proper Jacobian chain rule implementation, (3) Adaptive PID control with error-based damping, (4) Topology-driven strategy selection including progressive activation for series nonlinear circuits, (5) Comprehensive evaluation demonstrating robustness across six circuit categories. The framework is compatible with industry-standard IBIS models and provides a complete solution for modern circuit simulation challenges.
\end{abstract}

\begin{IEEEkeywords}
Circuit simulation, SPICE, nonlinear analysis, Newton-Raphson method, convergence, logarithmic transformation, topology analysis
\end{IEEEkeywords}

\section{Introduction}
\IEEEPARstart{M}{odern} electronic circuits present unprecedented challenges for simulation tools. The exponential growth in device complexity, combined with extreme parameter ranges in contemporary components, has pushed traditional SPICE-like simulators beyond their convergence limits. Light-emitting diodes (LEDs) exemplify this challenge, with saturation currents ($I_s$) now ranging from $10^{-12}$ A in older devices to $10^{-38}$ A or smaller in modern high-efficiency variants. Traditional Newton-Raphson-based solvers fail catastrophically when faced with such extreme nonlinearities, unable to converge even with extensive manual tuning.

The fundamental challenge stems from the exponential nature of semiconductor device equations. The Shockley diode equation:
\begin{equation}
I = I_s(e^{V/nV_t} - 1)
\end{equation}
produces Jacobian matrix condition numbers exceeding $10^{15}$ for modern LEDs, causing numerical overflow and convergence failure. While various approaches have been proposed, including source stepping \cite{kundert1995}, pseudo-transient analysis \cite{dong2007}, and homotopy methods \cite{watson1989}, none have successfully handled the extreme parameter ranges found in contemporary circuits.

% Continue with rest of introduction...

\section{Related Work and Background}

\subsection{Traditional SPICE Convergence Methods}
The Newton-Raphson method has been the cornerstone of circuit simulation since SPICE's inception \cite{nagel1973}. For a system of nonlinear equations $\mathbf{F}(\mathbf{x}) = 0$, the method iteratively updates the solution:
\begin{equation}
\mathbf{x}_{k+1} = \mathbf{x}_k - \mathbf{J}^{-1}(\mathbf{x}_k)\mathbf{F}(\mathbf{x}_k)
\end{equation}
where $\mathbf{J}$ is the Jacobian matrix.

% Continue with related work...

\section{GLACIER: Gradient Logarithmic Adaptive Solver}

\subsection{Mathematical Foundation}
The core innovation of GLACIER is the integration of logarithmic transformation directly into the Newton-Raphson iteration. For circuit equations with exponential terms, we transform selected variables to logarithmic space:
\begin{equation}
y_i = \log(x_i) \text{ for selected variables}
\end{equation}

\subsection{Phase 0: Gradient-Aware Region Identification}
Before beginning the Newton-Raphson iteration, GLACIER performs gradient analysis to identify regions of sharp nonlinearity:

\begin{algorithm}
\caption{Gradient-Aware Region Identification}
\begin{algorithmic}[1]
\REQUIRE Circuit netlist, ramp values
\ENSURE Sharp regions for log transformation
\STATE $\mathbf{gradients} \leftarrow []$
\FOR{$v$ in $\mathbf{ramp\_values}$}
    \STATE $\mathbf{x} \leftarrow \text{compute\_operating\_point}(\text{circuit}, v)$
    \STATE $g \leftarrow \text{compute\_log\_gradient}(\mathbf{x})$
    \STATE $\mathbf{gradients}.\text{append}(g)$
\ENDFOR
\STATE $\mathbf{sharp\_regions} \leftarrow \text{detect\_sharp\_transitions}(\mathbf{gradients})$
\RETURN $\mathbf{sharp\_regions}$
\end{algorithmic}
\end{algorithm}

% Continue with GLACIER details...

\section{MAESTRO: Topology-Aware Strategy Orchestration}

\subsection{Circuit Topology Analysis}
MAESTRO begins by analyzing circuit structure to identify patterns that benefit from specialized solving strategies.

% Add topology analysis details...

\section{Combined Framework Architecture}

The GLACIER-MAESTRO framework integrates both components synergistically:

\begin{figure}[!t]
\centering
% \includegraphics[width=\columnwidth]{framework_architecture.pdf}
\caption{GLACIER-MAESTRO framework architecture showing the interaction between topology analysis, strategy selection, and numerical solving.}
\label{fig:architecture}
\end{figure}

% Continue with framework details...

\section{Experimental Evaluation}

\subsection{Benchmark Suite}
We evaluate GLACIER-MAESTRO on 52 challenging circuits across six categories:

\begin{table}[!t]
\caption{Circuit Categories and Test Cases}
\label{tab:categories}
\centering
\begin{tabular}{lcc}
\toprule
Category & Circuits & Key Challenge \\
\midrule
Series Nonlinear & 15 & LED chains, $I_s \in [10^{-38}, 10^{-12}]$ \\
Parallel Arrays & 8 & Current sharing \\
Power Converters & 10 & Switching behavior \\
Cascaded Amplifiers & 7 & High gain \\
Bridge Circuits & 6 & Multiple diodes \\
Protection Circuits & 6 & Sharp transitions \\
\bottomrule
\end{tabular}
\end{table}

\subsection{Overall Results}

Table~\ref{tab:overall} shows convergence rates across all 52 circuits:

\begin{table}[!t]
\caption{Overall Convergence Results}
\label{tab:overall}
\centering
\begin{tabular}{lcccc}
\toprule
Solver & Converged & Rate & Mean Iter. & Mean Time \\
\midrule
Newton-Raphson & 19/52 & 36.5\% & 127.3 & 12.4 ms \\
GLACIER & 32/52 & 61.5\% & 1,847.2 & 423.7 ms \\
MAESTRO & 48/52 & 92.3\% & 318.7 & 67.2 ms \\
GLACIER-MAESTRO & 52/52 & \textbf{100\%} & 287.4 & 58.3 ms \\
\bottomrule
\end{tabular}
\end{table}

% Continue with results...

\section{Discussion}

\subsection{Why Does GLACIER-MAESTRO Succeed?}
The framework succeeds through synergistic combination of numerical and topological insights:
\begin{enumerate}
\item GLACIER addresses numerical challenges through logarithmic transformation
\item MAESTRO exploits circuit structure for efficient solving
\item Combined approach covers all cases with 100\% reliability
\end{enumerate}

% Continue discussion...

\section{Conclusion}

This paper presented GLACIER-MAESTRO, a comprehensive framework for robust circuit simulation that achieves 100\% convergence on challenging nonlinear circuits. By combining GLACIER's logarithmic transformation and gradient analysis with MAESTRO's topology-aware strategy orchestration, we solve previously intractable problems including LED circuits with saturation currents as low as $10^{-38}$ A.

% Add acknowledgments
\section*{Acknowledgment}
We utilized Claude (Anthropic) to assist with implementation, experimental automation, and manuscript preparation. All algorithmic innovations, experimental design, and scientific conclusions are the work of the human authors, who take full responsibility for the technical content.

% Bibliography
\bibliographystyle{IEEEtran}
\bibliography{references}

% Author biographies
\begin{IEEEbiography}[{\includegraphics[width=1in,height=1.25in,clip,keepaspectratio]{author1.jpg}}]{Author One}
Biography text here.
\end{IEEEbiography}

\begin{IEEEbiography}[{\includegraphics[width=1in,height=1.25in,clip,keepaspectratio]{author2.jpg}}]{Author Two}
Biography text here.
\end{IEEEbiography}

\end{document}